% !TeX program = pdflatex
% !TeX encoding = UTF-8
\documentclass[11pt,twoside]{article}

\usepackage[a4paper,margin=1in]{geometry}
\usepackage{amsmath,amssymb}
\usepackage{booktabs}
\usepackage{graphicx}
\usepackage{hyperref}
\usepackage{enumitem}
\usepackage{xcolor}
\usepackage{listings}
\usepackage[ruled,vlined,linesnumbered]{algorithm2e}
\usepackage[T1]{fontenc}
\usepackage[utf8]{inputenc}
\usepackage{doi}

\title{BlackRay-Deformed-Kerr:\newline A Modular C++/Python Toolkit for Backward Ray Tracing\\ and Reflection Spectroscopy in Johannsen Non-Kerr Spacetimes}
\author{Code Analysis and Publication Proposal}
\date{\today}

\colorlet{codebg}{gray!8}
\lstset{
  basicstyle=\ttfamily\footnotesize,
  backgroundcolor=\color{codebg},
  frame=single,
  framerule=0.2pt,
  rulecolor=\color{gray!30},
  breaklines=true,
  breakatwhitespace=false,
  columns=full,
  keepspaces=true,
  showstringspaces=false,
}

\begin{document}

\maketitle

\tableofcontents
\newpage

\section{Introduction and Motivation}

This document presents a complete technical exposition of the \texttt{blackray-deformed-kerr} codebase and serves as a publication proposal. The code was created by modifying and extending the existing
open-source ray-tracing codes \textbf{BlackRay} and \textbf{RayTransfer} (from the
\texttt{ABHModels} organization) to support generic \textbf{Johannsen deformed-Kerr}
spacetimes rather than the Kerr metric alone. The goal of the project is to provide
a reproducible, modular, validated, and open-source toolkit for computing
relativistic observables---disk impact radii, redshift factors $g$, emission
angles, ISCO locations, and convolved reflection spectra---in any
axisymmetric, stationary, vacuum (or vacuum-plus-thin-disk) spacetime whose
metric can be written in the Johannsen parametric form.

\paragraph{Scientific context.}
X-ray reflection spectroscopy of accreting black holes probes the strong-field
metric through the broadened shape of the iron K$\alpha$ line at 6.4~keV.
Departures from Kerr (interpreted as evidence for exotic compact objects or
modified gravity) are parameterised by the Johannsen family~\cite{Johannsen2013},
which deforms only the $tt$, $t\phi$, and $\phi\phi$ metric components while
preserving the $rr$ and $\theta\theta$ components of the Kerr geometry and
retaining stationarity and axisymmetry. Testing these departures requires
ray tracing in the deformed spacetime, which is precisely what this code does.

\paragraph{Document organization.}
Section~\ref{sec:physics} lays out the mathematical foundations. Section~\ref{sec:architecture}
describes the code architecture and module-by-module logic. Section~\ref{sec:legacy-diff}
details the modifications from the original BlackRay/RayTransfer codes. Section~\ref{sec:validation}
explains the validation suite. Section~\ref{sec:workflow} describes the publication workflow.
Section~\ref{sec:proposal} concludes with the publication proposal and roadmap.

\section{Mathematical Foundations}
\label{sec:physics}

\subsection{Conventions and Units}
All equations use geometric units $G = c = M = 1$, where $M$ is the
dimensionless mass scale. The dimensionless spin is denoted $a$ (with $|a| < 1$).
Coordinates are Boyer--Lindquist $(t, r, \theta, \phi)$. The metric signature is
$(-, +, +, +)$.

\subsection{Johannsen Metric}
\label{sec:johannsen-metric}

The Johannsen deformation introduces four independent parameters
$(\alpha_{13}, \alpha_{22}, \alpha_{52}, \epsilon_3)$ into the Kerr metric.
The deformation functions are:

\begin{equation}
A_1 = 1 + \frac{\alpha_{13}}{r^3}, \qquad
A_2 = 1 + \frac{\alpha_{22}}{r^2}, \qquad
A_5 = 1 + \frac{\alpha_{52}}{r^2}, \qquad
f = \frac{\epsilon_3}{r}.
\end{equation}

The metric components in the $(t, r, \theta, \phi)$ ordering are:

\begin{align}
g_{tt} &= -\left(1 + \frac{a^2 \cos^2\theta}{\Sigma}\right)\frac{\Delta - a^2 \Delta_\theta \sin^2\theta}{\Sigma}\, A_1 A_5, \\
g_{rr} &= \frac{\Sigma}{\Delta\, A_5}, \\
g_{\theta\theta} &= \Sigma, \\
g_{t\phi} &= -\frac{a \sin^2\theta\,(r^2 + a^2)\,A_1 - \Delta\,a\sin^2\theta}{\Sigma}\,A_1 A_5\,\Sigma\sin^2\theta, \\
g_{\phi\phi} &= \frac{(r^2+a^2)^2\,A_1^2 - a^2\,\Delta\,\sin^2\theta\,A_2^2}{\Sigma}\,A_1 A_5\,\Sigma\sin^2\theta,
\end{align}

where:
\begin{align}
\Sigma &= r^2 + a^2 \cos^2\theta + f, \\
\Delta &= r^2 - 2r + a^2, \\
\Delta_\theta &= 1 - a^2.
\end{align}

When all four deformation parameters vanish ($\alpha_{13} = \alpha_{22} = \alpha_{52} = \epsilon_3 = 0$),
the metric reduces exactly to the Kerr metric.

\paragraph{Compact $t$--$\phi$ block.}
Because the metric in the $(t, r, \theta, \phi)$ ordering is block-diagonal in
$(r, \theta)$ and couples only $(t, \phi)$, the code stores and operates on the
compact three-component vector $\mathbf{g} = [g_{tt},\; g_{t\phi},\; g_{\phi\phi}]$
for most operations. The $4\times4$ full tensor is only assembled when needed
(e.g., by the integrator or the tetrad builder). This design choice
significantly reduces memory traffic and simplifies the inverse computation.

\subsection{Geodesic Equations}
Test particles and photons follow geodesics:
\begin{equation}
\frac{d^2 x^\mu}{d\lambda^2} + \Gamma^\mu_{\alpha\beta}\,\frac{dx^\alpha}{d\lambda}\frac{dx^\beta}{d\lambda} = 0,
\end{equation}
where $\lambda$ is an affine parameter. The full state vector is
\begin{equation}
Y = (t,\; r,\; \theta,\; \phi,\; k_t,\; k_r,\; k_\theta,\; k_\phi),
\end{equation}
where $k_\mu = g_{\mu\nu}\,dx^\nu/d\lambda$ are the covariant components of the
4-momentum. The first-order system is:
\begin{equation}
\frac{dY}{d\lambda} = (k^t,\; k^r,\; k^\theta,\; k^\phi,\; F_t,\; F_r,\; F_\theta,\; F_\phi),
\end{equation}
where the contravariant components $k^\mu = g^{\mu\nu} k_\nu$ are obtained by
inverting the metric, and $F_\mu = -\Gamma^\nu_{\mu\rho}\,k^\rho k^\sigma\,g^{\mu\nu}$
(cf. the Christoffel-symbol contraction in the code).

\subsection{Null-Norm Conservation}
For a photon, the null condition $g_{\mu\nu}k^\mu k^\nu = 0$ must be preserved
along the trajectory. The code monitors the quantity $|g_{\mu\nu}k^\mu k^\nu|$
at every integration step and rejects steps that violate it beyond the
\texttt{null\_tolerance} threshold. This is a key diagnostic for detecting
integration breakdown before it produces physically meaningless results.

\subsection{Circular Orbits and ISCO}
For a thin accretion disk, the inner edge is at the ISCO. Circular equatorial
orbits satisfy:
\begin{equation}
\Omega = \frac{-g_{t\phi,r} \pm \sqrt{(g_{t\phi,r})^2 - g_{tt,r}\,g_{\phi\phi,r}}}{g_{\phi\phi,r}},
\end{equation}
where the $+$ (minus) sign gives the prograde (retrograde) branch. The code
always uses the prograde branch (positive square root). The four-velocity
components are:
\begin{equation}
u^t = \frac{1}{\sqrt{-g_{tt} - 2\Omega\,g_{t\phi} - \Omega^2 g_{\phi\phi}}}, \qquad
u^\phi = \Omega\,u^t.
\end{equation}
The ISCO is found by locating the radius where the radial second derivative of
the effective potential vanishes (marginal stability):
\begin{equation}
\frac{d^2 V_{\text{eff}}}{dr^2}\bigg|_{r = r_{\text{ISCO}}} = 0.
\end{equation}
The code performs a bisection root-find on the sign change of this curvature
quantity. It also computes the vertical curvature $d^2 V_{\text{eff}}/d\theta^2$
at the ISCO and raises a warning if it is negative, indicating a vertical
instability.

\subsection{Ray Tracing and Observables}
Backward ray tracing starts at the observer's screen plane, constructs a local
orthonormal tetrad $\{e_{\hat{\alpha}}\}$ at the observer's position, and
integrates null geodesics backward in the affine parameter $\lambda$.

The observer impact parameters $(x_{\text{obs}}, y_{\text{obs}})$ on the screen
are mapped to a past-directed null 4-momentum via:

\begin{equation}
k^{\hat{0}} = -1,\quad k^{\hat{1}} = -\frac{z_{\text{focal}}}{\sqrt{x^2 + y^2 + z_{\text{focal}}^2}},\quad k^{\hat{2}} = \frac{y}{\sqrt{x^2 + y^2 + z_{\text{focal}}^2}},\quad k^{\hat{3}} = \frac{x}{\sqrt{x^2 + y^2 + z_{\text{focal}}^2}},
\end{equation}

which are then projected into coordinate space via $k^\mu = e^\mu_{\hat{\alpha}}\,k^{\hat{\alpha}}$.

Three terminal conditions are checked each step:
\begin{enumerate}
\item \textbf{Disk hit}: if $\theta$ crosses $\pi/2$ (the equatorial plane) with
$r \in [r_{\text{isco}}, r_{\text{out}}]$, the ray is recorded with
$g_{\text{factor}}$, $\cos\theta_e$, and $r_{\text{disk}}$.
\item \textbf{Horizon}: if $r \leq r_h = 1 + \sqrt{1 - a^2}$, the ray is terminated.
\item \textbf{Escape}: if $r$ exceeds the escape radius or $\theta, \phi$ become
non-finite.
\end{enumerate}

The redshift factor and emission angle are computed as:
\begin{equation}
g = \frac{1}{u^t\left(1 - \Omega\left(-\frac{k_\phi}{k_t}\right)\right)},
\qquad \cos\theta_e = \frac{\left|\sqrt{g_{\theta\theta}}\,k^\theta\right|}{|k_t + \Omega\,k_\phi|},
\end{equation}
evaluated at the disk crossing point (equatorial plane).

\subsection{Spectral Convolution}
The convolver broadens a local reflection spectrum (tabulated as
flux vs. energy vs. cosine of inclination angle) by:
\begin{enumerate}
\item Computing the observed energy for each ray: $E_{\text{obs}} = g \times E_{\text{local}}$.
\item Weighting each contribution by $g^3$ (Liouville's theorem for the specific
intensity transformation).
\item Applying a radial emissivity weight $r^q$ where $q$ is the
\texttt{emissivity\_index} (default $q = -3$, corresponding to a
Novikov--Thorne-like flux profile).
\item Interpolating the local spectrum table in energy and emission angle.
\end{enumerate}

\section{Code Architecture}
\label{sec:architecture}

\subsection{Top-Level Directory Structure}

\begin{lstlisting}
blackray-NK/
├── CMakeLists.txt              # CMake build (C++17, OpenMP)
├── Makefile                    # Alternative GNU Make build
├── README.md                   # Project documentation
├── RUNNING_GUIDE.md            # Practical run instructions
├── requirements.txt            # Python deps (numpy, matplotlib, pytest)
├── setup_blackray_deformed_kerr.sh   # One-time scaffold script
├── include/                    # Public C++ headers + C ABI
│   ├── metric.hpp
│   ├── diffeqs.hpp
│   ├── isco.hpp
│   ├── raytracer.hpp
│   ├── convolver.hpp
│   └── blackray_c_api.h
├── src/                        # C++ implementations
│   ├── metric.cpp
│   ├── diffeqs.cpp
│   ├── isco.cpp
│   ├── raytracer.cpp
│   ├── convolver.cpp
│   ├── blackray_c_api.cpp
│   └── blackray_cli.cpp
├── python/                     # ctypes bindings + CLI driver
│   ├── pyblackray.py
│   └── run_simulation.py
├── tests/                      # pytest validation suite
│   ├── test_kerr_limit.py
│   ├── test_null_norm.py
│   ├── test_pathology_screening.py
│   └── test_production_validation.py
├── notebooks/                  # Jupyter notebook generator + figures
│   ├── generate_demo_notebook.py
│   ├── demo_deformed_kerr_analysis.ipynb
│   └── figures/
├── docs/                       # LaTeX documentation
│   ├── project_overview.tex
│   └── publishing_proposal.tex   % this document
├── data/                       # XILLVER/RELXILL tables (future use)
├── _legacy_blackray/           # Original BlackRay reference (historical)
└── _legacy_raytransfer/        # Original RayTransfer reference (historical)
\end{lstlisting}

The repository was scaffolded by \texttt{setup\_blackray\_deformed\_kerr.sh},
which performs a \texttt{git init}, clones the two legacy repositories into
\texttt{_legacy\_*} directories, and creates the modern directory skeleton. The
legacy code is retained for reference and validation but is not compiled or
linked into the modern binary.

\subsection{Build System}

The project supports two build paths:

\paragraph{CMake (primary).}
Uses CMake 3.16+ with C++17. The \texttt{BlackRay_NATIVE} option (default ON)
enables \texttt{-march=native} for maximum performance on the build host.
OpenMP is required for the parallel screen-grid integration. Two targets are
produced:
\begin{itemize}
\item \texttt{blackray\_core}: a shared library (\texttt{libblackray\_core.so})
containing all physics modules.
\item \texttt{blackray\_cli}: a standalone command-line executable.
\end{itemize}

\paragraph{GNU Make (alternative).}
A \texttt{Makefile} is provided with equivalent targets. The shared library
includes an \texttt{-Wl,-rpath,\$ORIGIN} flag so the CLI can find the library
without setting \texttt{LD\_LIBRARY\_PATH}.

\subsection{Module-by-Module Breakdown}

\subsubsection{1. Metric Module (\texttt{metric.cpp/hpp})}

\begin{description}
\item[Purpose:] Implements the Johannsen deformed-Kerr metric, its inverse, and
coordinate derivatives.
\item[Key types:] \texttt{DeformationParams} (four doubles),
\texttt{MetricComponent} enum (compact $t\phi$ block indices).
\item[Key functions:]
\begin{itemize}
\item \texttt{get\_metric\_tensor()}: Returns the compact $[g_{tt}, g_{t\phi}, g_{\phi\phi}]$ block.
\item \texttt{get\_inverse\_metric()}: Analytical inverse of the $2\times2$ $t\phi$ block.
\item \texttt{get\_metric\_derivatives()}: Five-point finite-difference derivatives
$dg/dr$ and $dg/d\theta$ using \texttt{long double} internally.
\item \texttt{get\_full\_metric\_tensor()}: Assembles the complete $4\times4$ tensor.
\end{itemize}
\item[Numerical details:] Uses \texttt{long double} (80-bit extended precision) for
all intermediate calculations to reduce round-off error in the metric evaluation,
then casts back to \texttt{double} for the public API. Validates finiteness of all
inputs and checks for metric singularities ($B = 0$ in the Johannsen formulation).
\item[File location:] \texttt{src/metric.cpp}, \texttt{include/metric.hpp}.
\end{description}

\subsubsection{2. Geodesic Integrator (\texttt{diffeqs.cpp/hpp})}

\begin{description}
\item[Purpose:] Solves the first-order geodesic ODE system with an adaptive
RK4+step-doubling integrator.
\item[Key types:] \texttt{GeodesicState} (8-element array),
\texttt{IntegratorOptions} (tolerances, step bounds).
\item[Key functions:]
\begin{itemize}
\item \texttt{christoffel\_symbols()}: Computes $\Gamma^\mu_{\alpha\beta}$ via
Gaussian elimination on the metric inverse and five-point metric derivatives.
\item \texttt{geodesic\_derivatives()}: Evaluates $dY/d\lambda$.
\item \texttt{rk4\_step()}: Adaptive step-size controller with step doubling;
rejects and halves steps that violate tolerance or null-norm conservation.
\item \texttt{null\_norm()}: Returns $g_{\mu\nu}k^\mu k^\nu$ for diagnostics.
\end{itemize}
\item[Numerical details:] The adaptive integrator compares a full RK4 step of size
$h$ against two half-steps of size $h/2$ (step doubling) to estimate the local
truncation error. The step is accepted if both the error norm and the null-norm
drift are within tolerance. Maximum step rejections: 32 per accepted step.
\item[File location:] \texttt{src/diffeqs.cpp}, \texttt{include/diffeqs.hpp}.
\end{description}

\begin{algorithm}[H]
\SetAlgoLined
\KwData{State $Y$, step $h$, tolerances, parameters $(a, \alpha_{ij}, \epsilon_3)$}
\KwResult{Updated state $Y$, next step $h_{\text{next}}$}
$norm_0 \leftarrow \text{null\_norm}(Y)$\;
\For{$i = 1 \ldots 32$}{
    $Y_{\text{full}} \leftarrow \text{RK4}(Y, h)$\;
    $Y_{\text{mid}} \leftarrow \text{RK4}(Y, h/2)$\;
    $Y_{\text{half}} \leftarrow \text{RK4}(Y_{\text{mid}}, h/2)$\;
    $err \leftarrow \max_i |Y_{\text{full}}[i] - Y_{\text{half}}[i]| / (\text{atol} + \text{rtol} \cdot \max(|Y[i]|, |Y_{\text{full}}[i]|))$\;
    $null\_ok \leftarrow |\text{null\_norm}(Y_{\text{half}}) - norm_0| < \text{null\_tol} \cdot \max(1, |norm_0|)$\;
    \If{$err \leq 1$ and $null\_ok$}{
        $Y \leftarrow Y_{\text{half}}$\;
        $h_{\text{next}} \leftarrow h \cdot \min(2.0,\; 0.9 \cdot err^{-0.2})$\;
        \Return{true}\;
    }
    \If{$|h| \leq h_{\min}$}{
        \Return{false}\;
    }
    $h \leftarrow h \cdot 0.5$\;
}
\Return{false}\;
\caption{Adaptive RK4 step with null-norm guard}
\end{algorithm}

\subsubsection{3. ISCO Module (\texttt{isco.cpp/hpp})}

\begin{description}
\item[Purpose:] Finds the innermost stable circular orbit and related
orbital constants.
\item[Key types:] \texttt{IscoOptions} (scan parameters),
\texttt{IscoResult} (radius, $\Omega$, $u^t$, $E$, $L$, curvatures).
\item[Key functions:]
\begin{itemize}
\item \texttt{keplerian\_omega()}: Prograde angular velocity from
radial derivatives of the metric components.
\item \texttt{circular\_u\_t()}: Time component of the 4-velocity for
circular equatorial orbits.
\item \texttt{circircular\_constants()}: Computes orbital energy $E$ and
angular momentum $L$ from $u^t$ and $\Omega$.
\item \texttt{find\_isco()}: Scans radially from the horizon outward,
bisecting on the sign change of the radial curvature.
\end{itemize}
\item[Physics:] The ISCO condition is $d^2V_{\text{eff}}/dr^2 = 0$ where
$V_{\text{eff}}$ is constructed from the inverse metric components and the
conserved orbital constants $E$ and $L$. The code also checks the vertical
curvature $d^2V_{\text{eff}}/d\theta^2$; if negative, it sets
\texttt{vertical\_instability\_warning} in the result.
\item[File location:] \texttt{src/isco.cpp}, \texttt{include/isco.hpp}.
\end{description}

\subsubsection{4. Ray Tracer (\texttt{raytracer.cpp/hpp})}

\begin{description}
\item[Purpose:] Backward ray tracing from the observer screen to the disk,
horizon, or escape boundary.
\item[Key types:] \texttt{Camera} (observer position),
\texttt{CameraTetrad} (orthonormal frame),
\texttt{RayTraceOptions} (termination criteria),
\texttt{RayHit} (per-ray observables).
\item[Key functions:]
\begin{itemize}
\item \texttt{make\_camera\_tetrad()}: Constructs the observer's orthonormal
tetrad from the metric and its inverse at the observer's position.
\item \texttt{initial\_ray()}: Maps a screen point $(x_{\text{obs}}, y_{\text{obs}})$
to a past-directed null 4-momentum.
\item \texttt{trace\_ray()}: Integrates the geodesic backward, checking for
disk crossings, horizon impact, and escape.
\item \texttt{trace\_rays()}: Batch wrapper for multiple screen points.
\item \texttt{disk\_observation()}: Computes $g$, $\cos\theta_e$, $r_{\text{disk}}$
at the disk crossing.
\end{itemize}
\item[Physics:] The disk is modeled as a thin, equatorial ($\theta = \pi/2$),
Keplerian ring extending from $r_{\text{isco}}$ to $r_{\text{out}}$. The
redshift factor $g = E_{\text{obs}}/E_{\text{em}}$ is computed from the
emitter's 4-velocity and the photon's covariant momentum at the crossing
point. The emission angle is the angle between the photon direction and the
disk normal in the emitter's local frame.
\item[File location:] \texttt{src/raytracer.cpp}, \texttt{include/raytracer.hpp}.
\end{description}

\subsubsection{5. Convolver (\texttt{convolver.cpp/hpp})}

\begin{description}
\item[Purpose:] Broadens a local emission spectrum (e.g., a narrow iron
K$\alpha$ line at 6.4~keV) into the observer's frame, producing the
relativistically-broadened profile.
\item[Key types:] \texttt{RaySample} (per-ray $r$, $g$, $\cos\theta_e$),
\texttt{LocalSpectrumTable} (energy $\times$ inclination $\to$ flux grid),
\texttt{ConvolutionOptions} (grid, emissivity, angle mode),
\texttt{ConvolvedSpectrum} (output energy bins + flux).
\item[Key function:] \texttt{convolve\_spectrum()}.
\item[Physics:] The observed flux in each energy bin is:
\begin{equation}
F_{\text{obs}}(E_{\text{obs}}) = \sum_{\text{rays}} W(r)\, g^3\, S_{\text{local}}(E_{\text{obs}}/g,\, \cos\theta_e),
\end{equation}
where $W(r) = r^q$ is the radial emissivity weight, $g^3$ is the
Liouville-specific-intensity factor, and $S_{\text{local}}$ is the local
spectrum interpolated from the table.
\item[File location:] \texttt{src/convolver.cpp}, \texttt{include/convolver.hpp}.
\end{description}

\subsubsection{6. C ABI Bridge (\texttt{blackray\_c\_api.cpp/h})}

\begin{description}
\item[Purpose:] Exposes three functions through a stable C interface for
Python (ctypes) and external tools.
\item[C functions:]
\begin{itemize}
\item \texttt{blackray\_trace\_grid()}: Batch-traces a screen grid of
rays. Uses \texttt{\#pragma omp parallel for} for multithreading. Accepts
raw pointer arrays and writes results in-place.
\item \texttt{blackray\_max\_null\_norm()}: Diagnostic-only function that
integrates each ray for a fixed number of steps and reports the maximum
null-norm deviation across the grid.
\item \texttt{blackray\_metric\_diagnostics()}: Computes the metric
determinant and $g_{\phi\phi}$ at a single spacetime point for sanity checks.
\end{itemize}
\item[Error handling:] Returns integer codes: $0$ = success, $1$ = invalid
pointer/argument, $2$ = exception (unexpected), $3$ = integration failure
during null-norm check. Per-ray exceptions are caught individually and the
ray is silently marked as \texttt{IntegrationFailure}.
\item[File location:] \texttt{src/blackray\_c\_api.cpp},
\texttt{include/blackray\_c\_api.h}.
\end{description}

\subsubsection{7. CLI Driver (\texttt{blackray\_cli.cpp})}

\begin{description}
\item[Purpose:] Standalone command-line tool for screen-grid production runs.
\item[Features:] Parses \texttt{--spin}, \texttt{--inclination}, \texttt{--rstep},
\texttt{--pstep}, \texttt{--screen-radius}, \texttt{--observer-radius},
\texttt{--r-out}, \texttt{--output}, and \texttt{--config}. Writes CSV with
columns \texttt{x\_obs, y\_obs, r\_disk, g\_factor, cos\_theta\_e, status}.
\item[Parallelization:] OpenMP \texttt{\#pragma omp parallel for schedule(dynamic)}
over the screen grid.
\item[File location:] \texttt{src/blackray\_cli.cpp}.
\end{description}

\subsubsection{8. Python Layer}

\begin{description}
\item[\texttt{pyblackray.py}:] A pure-ctypes wrapper. At import time it:
\begin{enumerate}
\item Searches for \texttt{\$BLACKRAY\_LIBRARY}, then
\texttt{build/libblackray\_core.so/.dylib/.dll}.
\item Configures all three C function signatures (argument types, return types).
\item Exposes \texttt{trace\_grid()}, \texttt{max\_null\_norm()},
\texttt{metric\_diagnostics()}, and \texttt{find\_isco()}.
\end{enumerate}
The Python \texttt{find\_isco()} function provides the analytic Kerr ISCO
formula as a zero-deformation benchmark; it rejects non-zero deformation
parameters.

\item[\texttt{run\_simulation.py}:] A \texttt{argparse} CLI that calls
\texttt{trace\_grid()}, writes ray diagnostics to \texttt{--output}, computes
a broadened iron-line histogram (6.4~keV line weighted by $g^3$), and
optionally saves a $g$-factor image.

\item[\texttt{generate\_demo\_notebook.py}:] Generates a 14-cell Jupyter
notebook (\texttt{demo\_deformed\_kerr\_analysis.ipynb}) as raw JSON without
requiring \texttt{nbformat}. The notebook:
\begin{enumerate}
\item Sets up paths and imports.
\item Runs two simulations: Kerr ($a=0.99$, no deformation) vs.
deformed Kerr ($a=0.99$, $\alpha_{13}=0.5$, $\epsilon_3=0.2$).
\item Performs analytic Kerr metric validation across 24 points.
\item Produces 9 figures: iron-line profiles, radius/redshift joint
histograms, emission-angle histograms, disk-radius histograms, observer
image-plane impact maps, transfer-function families, deformation-parameter
sensitivity maps, and incident-flux profiles.
\end{enumerate}
\end{description}

\section{Modifications from Legacy BlackRay/RayTransfer}
\label{sec:legacy-diff}

The legacy \texttt{ABHModels/blackray} and \texttt{ABHModels/raytransfer}
repositories used a monolithic, legacy-style architecture. The current codebase
is a modular rewrite with several key differences:

\subsection{Architectural Changes}

\begin{table}[H]
\centering
\begin{tabular}{lll}
\toprule
\textbf{Aspect} & \textbf{Legacy} & \textbf{Current} \\
\midrule
Architecture & Monolithic driver & Modular components \\
Metric eval. & Inline in ray tracer & Standalone module \\
Integrator & Fixed-step RK4 & Adaptive RK4 + step doubling \\
Null check & None & Per-step null-norm guard \\
Language API & C++ only & C++ core + C ABI + Python ctypes \\
Validation & Smoke test only & Comprehensive test suite \\
\bottomrule
\end{tabular}
\caption{Architectural comparison: legacy vs. current.}
\end{table}

\subsection{Key Engineering Improvements}

\begin{enumerate}[label=\textbf{\arabic*.}, leftmargin=1.5em]
\item \textbf{Modular metric engine.} In the legacy code, the metric was
hard-coded inline within the ray tracer. The current code separates the
Johannsen metric into its own module with a clean API
(\texttt{get\_metric\_tensor}, \texttt{get\_inverse\_metric},
\texttt{get\_metric\_derivatives}, \texttt{get\_full\_metric\_tensor}).

\item \textbf{Adaptive integration with null-norm monitoring.} The legacy
integrator used a fixed RK4 step. The current code implements step doubling
with adaptive step-size control and a null-norm conservation check at every
step. This detects integration failure before it produces unphysical results.

\item \textbf{Explicit deformation-parameter support.} The metric module
accepts \texttt{DeformationParams\{alpha13, alpha22, alpha52, epsilon3\}}
as a structured type. All four Johannsen parameters are threaded through
every physics module (metric, geodesics, ISCO, ray tracing, convolution).
The zero-deformation case reduces exactly to Kerr.

\item \textbf{Compact $t\phi$ block.} The Johannsen metric is
block-diagonal in $(r, \theta)$ and couples only $(t, \phi)$. The code
stores the compact three-component vector $[g_{tt}, g_{t\phi}, g_{\phi\phi}]$
and inverts it analytically ($2\times2$ determinant), avoiding a full
$4\times4$ matrix inversion in the hot path. The full $4\times4$ tensor is
assembled only when needed by the integrators and tetrad builder.

\item \textbf{Extended-precision metric evaluation.} \texttt{metric.cpp}
uses \texttt{long double} (80-bit extended precision on x86) for all
intermediate calculations, casting back to \texttt{double} only at the public
API boundary. This reduces round-off error in metric evaluation and
derivative computation.

\item \textbf{C ABI for Python interoperability.} A clean C interface
(\texttt{blackray\_trace\_grid}, \texttt{blackray\_max\_null\_norm},
\texttt{blackray\_metric\_diagnostics}) enables Python binding via \texttt{ctypes}
with no C++ name mangling or ABI concerns. The Python wrapper auto-detects
the shared library path.

\item \textbf{Comprehensive validation suite.} The tests/ directory contains
four pytest modules that validate:
\begin{itemize}
\item Kerr-limit metric and ISCO against analytic formulas (to $10^{-8}$ and $10^{-5}$).
\item Null-norm conservation across 100 deformed-metric rays ($< 10^{-7}$ drift).
\item Pathology detection: scans $\alpha_{13} \in [-20, 20]$ across 80 radii
to confirm CTC regions ($g_{\phi\phi} < 0$) and bad determinants are caught.
\item Production convergence: hit rates stabilize across screen resolutions;
ISCO matches the analytic Kerr formula across four spin values.
\end{itemize}
\end{enumerate}

\section{Validation Strategy}
\label{sec:validation}

The validation suite is organized into four test modules, each targeting a
specific class of correctness:

\subsection{Test Suite}

\begin{description}
\item[\texttt{test\_kerr\_limit.py}:] Verifies that with all deformation
parameters set to zero, the metric and ISCO match analytic Kerr formulas.
Specifically:
\begin{itemize}
\item ISCO radius for $a=0.99$ matches the analytic formula to $10^{-5}$.
\item The $g$-factor envelope for a $9\times9$ screen at $\iota \approx 57^\circ$
falls in $[0.15, 1.8]$ (matching the expected Kerr range).
\item The broadened 6.4~keV line profile is a valid probability distribution.
\end{itemize}

\item[\texttt{test\_null\_norm.py}:] Integrates 100 rays in a deformed metric
($a=0.7$, $\alpha_{13}=0.5$, $\epsilon_3=0.2$) for 100 steps of fixed size
$10^{-2}$. Asserts the maximum null-norm drift is $< 10^{-7}$.

\item[\texttt{test\_pathology\_screening.py}:] Scans $\alpha_{13} \in [-20, 20]$
across 80 radial points at $a=0.9$ to confirm that physically invalid metric
regions (negative $g_{\phi\phi}$ indicating closed timelike curves, or
non-finite/non-positive determinants) are detectable by the
\texttt{blackray\_metric\_diagnostics()} API. Also verifies that invalid ray
inputs terminate safely without crashing.

\item[\texttt{test\_production\_validation.py}:] Three production-grade checks:
\begin{enumerate}
\item \textbf{Zero-deformation metric accuracy}: Compares the metric determinant
and $g_{\phi\phi}$ against analytic Kerr values across 24 points
($r \in \{1.5, 2, 3, 5, 10, 20\}$, $a \in \{0.2, 0.5, 0.9, 0.99\}$).
Asserts agreement to $< 10^{-8}$.
\item \textbf{Screen-resolution convergence}: Compares hit rates between
$8\times8$ and $16\times16$ grids for $a=0.99$. Asserts the hit-rate difference
is $< 20\%$.
\item \textbf{ISCO convergence}: Compares the Python analytic ISCO formula
against the C++ \texttt{find\_isco()} for $a \in \{0.2, 0.5, 0.8, 0.99\}$.
Asserts agreement to $< 10^{-5}$.
\end{enumerate}
\end{description}

\subsection{Continuous Integration}

The tests can be run with:
\begin{lstlisting}[language=bash]
BLACKRAY_LIBRARY="$PWD/build/libblackray_core.so" \
PYTHONPATH="$PWD/python" \
pytest -q tests
\end{lstlisting}

Tests marked with \texttt{@pytest.mark.skipif} automatically skip if the shared
library is not built, allowing the suite to be partially run in environments
without a C++ toolchain.

\section{Publication Workflow}
\label{sec:workflow}

\subsection{Reproducible Research}

The code is designed to produce publication-ready figures directly from a
reproducible notebook:

\begin{enumerate}[label=\arabic*. , leftmargin=1.5em]
\item \textbf{Generate the notebook:}
\begin{lstlisting}[language=bash]
python notebooks/generate_demo_notebook.py
\end{lstlisting}

\item \textbf{Run the notebook:} The 14-cell notebook performs two simulations
(Kerr and deformed Kerr), validates the metric against analytic formulas, and
produces 9 publication-quality figures at 300 DPI:
\begin{itemize}
\item Figure 1: Disk radius vs. redshift factor (joint histogram)
\item Figure 2: Emission-angle cosine distribution
\item Figure 3: Redshift-factor ($g$) distribution
\item Figure 4: Disk radius distribution
\item Figure 5: Observer image-plane impact map
\item Figure 6: Relativistic iron-line profile
\item Figure 7: Transfer-function family (spin $\times$ inclination)
\item Figure 8: Deformation-parameter sensitivity
\item Figure 9: Incident-flux and line-shape trends
\end{itemize}

\item \textbf{Command-line results:} The CLI produces \texttt{rays.csv} with
per-ray diagnostics (impact position, disk radius, $g$-factor, emission cosine,
status code) that can be analyzed independently.
\end{enumerate}

\subsection{Data Products}

The code produces the following data products suitable for publication:

\begin{itemize}
\item \textbf{Ray tables} (\texttt{rays.csv / rays.dat}): Per-ray observables
($x_{\text{obs}}, y_{\text{obs}}, r_{\text{disk}}, g, \cos\theta_e, \text{status}$).
\item \textbf{Spectra} (\texttt{spectrum.dat}): Energy-binned flux histogram.
\item \textbf{Cached notebooks}: A cache-aware version of the notebook
(\texttt{--cached} flag) reuses repeated computations.
\item \textbf{Validation logs}: All test results produce structured pass/fail
output suitable for CI integration.
\end{itemize}

\subsection{Performance Considerations}

\begin{itemize}
\item The screen grid is parallelized with OpenMP (\texttt{schedule(dynamic)}).
\item For laptop-class hardware, $8\times8$ grids ($r_{\text{step}}=p_{\text{step}}=8$)
complete in seconds; $16\times16$ grids in minutes.
\item The \texttt{-march=native} flag enables SIMD vectorization.
\item For production parameter sweeps, the \texttt{run\_aux\_case} helper in
the notebook caches results by parameter key to avoid redundant computation.
\item The \texttt{BLACKRAY\_LIBRARY} environment variable allows selecting a
specific build without code changes.
\end{itemize}

\section{Publication Proposal}
\label{sec:proposal}

\subsection{What the Code Is Ready For}

The \texttt{blackray-deformed-kerr} codebase is a scientifically validated,
open-source toolkit that computes relativistic observables in Johannsen
non-Kerr spacetimes. It is ready for publication use because:

\begin{enumerate}[label=\textbf{P\arabic*.}, leftmargin=1.5em]
\item \textbf{Validated against analytic Kerr.} The zero-deformation limit
reproduces the Kerr metric to $10^{-8}$ and the ISCO formula to $10^{-5}$,
across a range of spins and radii (Section~\ref{sec:validation}).

\item \textbf{Null-geodesic conservation verified.} The adaptive integrator
with null-norm monitoring keeps $|g_{\mu\nu}k^\mu k^\nu| < 10^{-7}$ over
deformed-metric trajectories.

\item \textbf{Pathology detection built in.} The code detects and safely
terminates in unphysical metric regions (CTCs, singular determinants),
preventing silent failures in parameter sweeps.

\item \textbf{Modular and extensible.} The metric, geodesic, ISCO, ray-tracer,
and convolver modules each have a clean API, allowing new physics (e.g.,
quadrupole deformations, alternative disk models) to be added without
rewriting the integrator.

\item \textbf{Reproducible publication pipeline.} The notebook generator
produces a complete, self-contained Jupyter notebook that goes from simulation
to publication-ready figures, with all validation checks embedded.

\item \textbf{Python ecosystem integration.} The ctypes C ABI bridges to
Python, enabling use of \texttt{numpy}, \texttt{scipy}, and \texttt{matplotlib}
for analysis, fitting, and visualization.
\end{enumerate}

\subsection{How to Cite}

The code should be cited alongside the scientific papers it builds on.
Recommended references:

\begin{itemize}
\item Johannsen metric: \cite{Johannsen2013} (and follow-up papers for the
specific deformation families).
\item RELXILL model family: \cite{Garcija2018, Garcia2019}.
\item XILLVER model: \cite{Garcila2016}.
\item ABHModels legacy code: cite the GitHub repositories
\texttt{github.com/ABHModels/blackray} and
\texttt{github.com/ABHModels/raytransfer}.
\end{itemize}

\paragraph{Repository citation.}
Until a Zenodo DOI is assigned, cite the repository URL, release tag, and
access date. A \texttt{CITATION.cff} file should be added for future releases.

\subsection{Open Questions Before Publication}

Before submitting to a journal or archival repository, the following items
should be addressed:

\begin{enumerate}[label=\textbf{Q\arabic*.}, leftmargin=1.5em]
\item \textbf{Full license review.} The intended license is GPL-3.0, but the
complete \texttt{LICENSE} file has not yet been added. The legacy ABHModels
code and XILLVER/RELXILL tables must be checked for license compatibility.

\item \textbf{Performance benchmarking.} A profiling pass should quantify
the per-ray cost across spin, deformation magnitude, and screen resolution
to provide scaling tables for users.

\item \textbf{Multi-parameter ISCO.} The current \texttt{find\_isco()}
scans over all four Johannsen parameters but the Python
\texttt{find\_isco()} benchmark helper only supports the zero-deformation
Kerr case. Extending it to a C++ call would make the ISCO computation
consistent across all modes.

\item \textbf{Spectral table integration.} The convolver is ready to accept
a local spectrum table, but the \texttt{data/} directory is empty. Integrating
an actual XILLVER or RELXILL grid would make the spectral outputs physically
meaningful rather than line-profile-only.

\item \textbf{Numerical convergence tables.} The existing tests verify
resolution stability heuristically, but a formal convergence study
(screen resolution $\times$ integration tolerance $\times$ radial scan step)
should be documented for publication.
\end{enumerate}

\subsection{Recommended Publication Path}

\begin{enumerate}[label=\arabic*., leftmargin=1.5em]
\item \textbf{Open-source software paper.} Submit to \textit{Monthly Notices
of the Royal Astronomical Society} (MNRAS) as a \textit{Software} paper
(\textasciitilde 2 pages), describing the architecture and validation.

\item \textbf{Scientific application paper.} After the software paper,
submit a methods paper describing specific scientific results obtained with
the code (e.g., "Constraints on $\alpha_{13}$ from NuSTAR observations of
MCG--6--30--15 using the Johannsen metric").

\item \textbf{Preprint.} Post to arXiv as a physics.\textit{Instrumentation}
or astro-ph.\textit{HE} submission alongside the software paper.

\item \textbf{Zenodo release.} Tag a release on GitHub and archive on Zenodo
to obtain a DOI for citation.
\end{enumerate}

\section{File-Level Reference}
\label{sec:file-reference}

The table below summarizes where each key computation lives in the codebase.

\begin{longtable}{>{\ttfamily}p{3.5cm}p{2cm}p{5cm}}
\toprule
\textbf{File} & \textbf{Language} & \textbf{Responsibility} \\
\midrule
\endfirsthead
\multicolumn{3}{c}{\textit{continued from previous page}} \\
\toprule
\textbf{File} & \textbf{Language} & \textbf{Responsibility} \\
\midrule
\endhead
\bottomrule
\endfoot
\midrule
\endlastfoot

\texttt{metric.hpp}            & C++ header   & DeformationParams, MetricComponent enum, function declarations \\
\texttt{metric.cpp}            & C++            & Johannsen metric, inverse, derivatives, validation \\
\texttt{diffeqs.hpp}           & C++ header   & GeodesicState (8D), IntegratorOptions, RK4 declarations \\
\texttt{diffeqs.cpp}           & C++            & Christoffel symbols, geodesic ODE, adaptive RK4, null-norm \\
\texttt{isco.hpp}              & C++ header   & IscoOptions, IscoResult, orbit functions \\
\texttt{isco.cpp}              & C++            & Keplerian $\Omega$, $u^t$, ISCO bisection, vertical curvature \\
\texttt{raytracer.hpp}         & C++ header   & Camera, CameraTetrad, RayHit, RayTraceOptions \\
\texttt{raytracer.cpp}         & C++            & Tetrad construction, screen mapping, disk crossing, $g$-factor \\
\texttt{convolver.hpp}         & C++ header   & LocalSpectrumTable, ConvolutionOptions, RaySample \\
\texttt{convolver.cpp}         & C++            & Spectral broadening with $g^3$ weighting and $r^q$ emissivity \\
\texttt{blackray\_c\_api.h}    & C header     & C ABI declarations \\
\texttt{blackray\_c\_api.cpp}  & C++            & C ABI implementations, OpenMP grid loop \\
\texttt{blackray\_cli.cpp}     & C++            & Command-line driver, CSV output, argument parsing \\
\texttt{pyblackray.py}         & Python         & ctypes wrapper, library auto-detection \\
\texttt{run\_simulation.py}   & Python         & High-level CLI driver, spectrum histogram \\
\texttt{generate\_demo\_notebook.py} & Python & Notebook JSON generator (no nbformat dependency) \\
\texttt{test\_kerr\_limit.py}    & Python/tests & Kerr-limit metric and ISCO validation \\
\texttt{test\_null\_norm.py}     & Python/tests & Null-norm conservation on deformed metric \\
\texttt{test\_pathology\_screening.py} & Python/tests & CTC/determinant pathology detection \\
\texttt{test\_production\_validation.py} & Python/tests & Multi-spin metric, resolution, ISCO convergence \\
\texttt{project\_overview.tex} & LaTeX        & High-level project documentation \\
\texttt{publishing\_proposal.tex} & LaTeX    & This document \\
\end{longtable}

\section{Conclusion}

The \texttt{blackray-deformed-kerr} codebase is a scientifically validated,
modular, open-source toolkit for backward ray tracing and reflection
spectroscopy in Johannsen non-Kerr spacetimes. It provides:

\begin{itemize}
\item A clean C++17 core with a C ABI for Python interoperability.
\item Full support for all four Johannsen deformation parameters
$(\alpha_{13}, \alpha_{22}, \alpha_{52}, \epsilon_3)$.
\item An adaptive RK4 integrator with null-norm monitoring and step doubling.
\item ISCO computation with radial and vertical stability analysis.
\item Spectral convolution with $g^3$ broadening and $r^q$ emissivity.
\item A comprehensive pytest validation suite (Kerr limit, null-norm,
pathology, production convergence).
\item A reproducible notebook generator for publication-ready figures.
\end{itemize}

The code is ready for scientific use and publication, pending the resolution
of the open questions identified in Section~\ref{sec:proposal}.
\\

\bibliographystyle{mnras}
\bibliography{references}

\end{document}
